\chapter{What gets recorded}
\label{chap:record}

The session produces two things: image files, which the software names and
stores by itself, and written observations, which nothing records unless a
person writes them down. This chapter specifies the second, because it is the
part that has previously gone missing.

\section{The record sheet}

Every field is filled during the session, not reconstructed afterwards. A
field that could not be measured is marked \emph{not measured} with the
reason, which is criterion A4.1. An empty cell is indistinguishable from a
forgotten one.

\begin{longtable}{@{}p{0.52\textwidth} p{0.13\textwidth} p{0.27\textwidth}@{}}
\caption[Session record sheet]{The session record sheet. Transcribed into the
results chapter after the session.}
\label{tab:record} \\
\toprule
\textbf{Datum} & \textbf{Unit} & \textbf{Value} \\
\midrule
\endfirsthead
\toprule
\textbf{Datum} & \textbf{Unit} & \textbf{Value} \\
\midrule
\endhead
\multicolumn{3}{@{}l}{\textit{Venue and date}} \\
Room, date, start and end time & & \\
People present and their roles & & \\
\addlinespace
\multicolumn{3}{@{}l}{\textit{Fixture}} \\
Make, model, modifier, count & text & \\
Ballast or driver, and the mode selected & text & \\
Dimming method and PWM frequency, if LED & text, Hz & \\
Warm-up time allowed before measuring & min & \\
Distance and angle of the fixture to the structure and to the cameras & m, \textdegree & \\
\addlinespace
\multicolumn{3}{@{}l}{\textit{Photometry}} \\
Illuminance at the tether position & lx & \\
Uniformity grid: minimum and maximum, and the grid spacing used & lx, m & \\
Measured colour temperature & K & \\
Measured CRI, and TLCI if available & & \\
Penumbra width at 1~m & mm & \\
Stability over the session: minimum and maximum & lx & \\
Stray-light check result & pass/fail & \\
Flicker test result, per camera & pass/fail & \\
\addlinespace
\multicolumn{3}{@{}l}{\textit{Geometry of the scene}} \\
Tether physical length, tape measured & m & \\
Tether diameter, measured & mm & \\
Tether material and surface finish & text & \\
Attachment points on the structure, described or sketched & & \\
Tether approach angle relative to the camera baseline & \textdegree & \\
Distance from the camera baseline to the nearest and farthest tether point & m & \\
Working volume of the structure: width, height, depth & m & \\
\addlinespace
\multicolumn{3}{@{}l}{\textit{Optics and cameras}} \\
Focus distance set, and how it was judged & m, text & \\
Aperture set & f/ & \\
Both rings locked and marked, before the session & yes/no & \\
Both marks unmoved, after the session & yes/no & \\
Per camera: exposure, gain, white balance & ms, dB & \\
Capture resolution and sensor mode used & px & \\
\addlinespace
\multicolumn{3}{@{}l}{\textit{Calibration}} \\
Session directory name & text & \\
Accepted views per tilt bin & count $\times$ 4 & \\
Relative standard deviation of $f_x$, $f_y$ & \% & \\
Principal point offset from centre & px, \% & \\
Reported reprojection RMS, for the record only & px & \\
Session reported complete by the tool & yes/no & \\
\addlinespace
\multicolumn{3}{@{}l}{\textit{Captures}} \\
Capture sets obtained, against Table \ref{tab:capture-plan} & S1--S8 & \\
Total stereo pairs & count & \\
Inter-camera delay: median and worst & ms & \\
Anything that went wrong, in prose & & \\
\bottomrule
\end{longtable}

\section{Why the scene geometry fields are there}

The block headed \emph{Geometry of the scene} is not documentation for its
own sake. The Blender validation deliverable selected a real stereo pair as
its reference and then found it could not reproduce the scene, because the
tether's approach angle, distance and attachment points had never been
written down \cite{linku_blender_validation}. Those four fields exist so that
this session's captures do not become unusable for the same reason.

The photometry block serves the same deliverable from the other side: its
laboratory-condition light is currently a placeholder, and these measurements
are what replaces it.

\section{Files, naming and metadata}

The application writes each stereo pair to its own timestamped directory
under \texttt{data/\allowbreak captures/}, with per-pair metadata including the capture
resolution, the baseline and the inter-camera delay. That naming is
automatic and should not be changed by hand. What has to be added is the
mapping from directory to capture set.

\begin{itemize}
    \item Keep a plain text file in the session folder mapping each capture
    directory to its set S1 to S8 and to a one-line description of what was
    being done at the time. Written during the session, not after.

    \item Commit the calibration session directory and every capture
    directory, with their metadata, which is criterion A4.5. The repository
    already versions \texttt{data/\allowbreak captures/} and
    \texttt{data/\allowbreak calibration/}, so this is the default behaviour rather than
    an exception, but it has to actually be done.

    \item Photograph the physical setup: the fixture position, the structure
    with the tether strung, the camera support in place, and the ring marks.
    These go in \texttt{Images/} of this deliverable.
\end{itemize}

\section{What is deliberately not recorded}

The angular velocity of the hand motion. There is no instrumentation for it,
and writing down an estimate would create a number that later analysis would
be tempted to treat as a measurement. Section \ref{sec:why-by-hand} gives the
reasoning. The rolling-shutter deliverable is where that quantity gets
measured, with a stage that either commands or encodes the rate
\cite{linku_rollingshutter}.
